\documentclass[9pt]{beamer}

\usepackage[
  backend=biber,
  style=authoryear,
]{biblatex}
\addbibresource{literature2.bib}
%\usepackage{bibunits}
 %\addbibresource{literature.bib}
% \usepackage{algorithm,algorithmic}
\usepackage{amsmath}
% \usepackage{stmaryrd}
% \usepackage{longtable}
% \usepackage{proof,mathpartir}
% \usepackage{fix-cm}
% \usepackage{bm}
% \usepackage{listings}
% \usepackage{amssymb}
% \usepackage{colortbl}
 \usepackage{tcolorbox}
% \usepackage{dsfont}
% \usepackage{amsfonts}
% \usepackage{ulem,xcolor}
% \usepackage{xspace}
% \usepackage{eqparbox}
% \usepackage{siunitx,amsmath}
\definecolor{safetyred}{RGB}{200,0,0}
\definecolor{traceblue}{RGB}{0,80,200}
\definecolor{codegray}{rgb}{0.95,0.95,0.95}
\definecolor{commentgreen}{rgb}{0.0,0.5,0.0}
\definecolor{keywordblue}{rgb}{0.0,0.0,0.8}
\definecolor{alertred}{rgb}{0.8,0.0,0.0}
\input{macrossl.tex}

% % Code Listing Settings
% \lstset{
%   backgroundcolor=\color{codegray},
%   commentstyle=\color{commentgreen},
%   keywordstyle=\color{keywordblue}\bfseries,
%   basicstyle=\ttfamily\small,
%   breaklines=true,
%   columns=fullflexible,
%   frame=single,
%   showstringspaces=false,
%   captionpos=b,
%   tabsize=2
% }
% \lstdefinelanguage{YAML}{
%   keywords={true, false, null, y, n},
%   keywordstyle=\color{darkgray}\bfseries,
%   basicstyle=\ttfamily\footnotesize,
%   commentstyle=\color{teal}\itshape,
%   stringstyle=\color{blue},
%   morestring=[b]',
%   morestring=[b]",
%   morecomment=[l]{\#},
%   moredelim=[l][\color{orange}]{\&},
%   moredelim=[l][\color{magenta}]{*},
%   moredelim=**[il][\color{red}]{:},
%   literate={---}{{\textcolor{gray}{---}}}3
% }

%\addbibresource{literature.bib}
% \newcommand\mCite[1]{[\cite{#1}, \citetitle{#1}]}
% \newcommand{\itemEq}[1]{%
% 	\begingroup%
% 	\setlength{\abovedisplayskip}{0pt}%
% 	\setlength{\belowdisplayskip}{0pt}%
% 	\parbox[c]{\linewidth}{\begin{flalign}#1&&\end{flalign}}%
% 	\endgroup}
% \renewcommand\algorithmiccomment[1]{%
% 	\hfill\#\ \eqparbox{COMMENT}{#1}%
% }



% Encoding
\usepackage[T1]{fontenc}
% UTF8 input
\usepackage[utf8]{inputenc}
% Use german language
\usepackage[english]{babel}
\renewcommand{\arraystretch}{1.3}

% Use the given theme
\usetheme{HLISP}
%\renewcommand{\algorithmicforall}{\textbf{foreach}}
\title{Formal Modeling and Analysis of Metric-Timed Systems and Multi-Agent Normative Specifications}
\subtitle{PhD Defense – Doctor of Natural Sciences (Dr. rer. nat)}

\author{
  \footnotesize\textbf{Supervisor:}\\ Prof. Dr. Martin Leucker \\
  \textbf{Second Referee:}\\ Prof. Dr. Gordon J. Pace \\
  \textbf{Chair of the Examination Committee:}\\ Prof. Dr. Stefan Fischer
}
\date{29 April 2026}

\begin{document}

% --- Title ---
\begin{frame}
  \titlepage
\end{frame}

\begin{frame}{Outline}
	\tableofcontents
  \end{frame}
\section{Introduction and Context}
\begin{frame}{The Cost of Normative Compliance}
  \begin{columns}
    \begin{column}{0.5\textwidth}
      \textbf{Motivation:}
      \begin{itemize}
        \item Normative systems govern behavior via \textbf{ought-to-do rules}.
        \item Manual verification creates a heavy economic burden.
        \item Real-world rules involve strict \textbf{time constraints} and \textbf{distributed agents}.
      \end{itemize}
      
      \vspace{0.3cm}
      \textbf{The Problem:}
      \begin{itemize}
        \item Fragments of laws often \textbf{conflict} (ambiguity).
        \item Accountability fails when multi-agent collaboration breaks down.
      \end{itemize}
    \end{column}

    \begin{column}{0.5\textwidth}
      \centering
      \begin{tikzpicture}[scale=0.8, transform shape, node distance=1.8cm]
        % Manual Process
        \node (rulebook) [rectangle, draw, thick, minimum width=2.5cm, minimum height=1.2cm, align=center, fill=gray!10] {Manual\\Rulebook};
        \node (maze) [rectangle, draw, thick, below of=rulebook, minimum width=2.5cm, minimum height=1.2cm, align=center, fill=red!20] {Administrative\\Bottleneck};
        \draw [->, thick] (rulebook) -- (maze);

        % Digital Process
        \node (digital) [rectangle, draw, thick, right=2cm of rulebook, minimum width=2.5cm, minimum height=1.2cm, align=center, fill=blue!10] {Formal\\Rules};
        \node (flow) [rectangle, draw, thick, below of=digital, minimum width=2.5cm, minimum height=1.2cm, align=center, fill=green!20] {Automated\\Compliance};
        \draw [->, thick] (digital) -- (flow);

        \draw [dashed, thick, gray] (1.25,-0.6) -- (1.25,-2.8);
      \end{tikzpicture}
    \end{column}
  \end{columns}
\end{frame}

\begin{frame}{Research Objectives}
  \begin{tcolorbox}[colback=blue!5, colframe=blue!80, title=RQ1 / Chapter 3: Static Conflict Analysis]
    How can we formally detect, measure, and explain conflicts between time-constrained obligations and prohibitions?
  \end{tcolorbox}
  
  
  
  \begin{tcolorbox}[colback=orange!5, colframe=orange!80, title=RQ2 / Chapter 3: Conflict Resolution]
    Can we automatically repair a partially conflicting normative specification without altering its original intended behavior?
  \end{tcolorbox}


  
  \begin{tcolorbox}[colback=green!5, colframe=green!80, title=RQ3 / Chapter 4: Dynamic Blame Attribution]
    How do we formalize and quantify responsibility when collaborative execution of norms fails in a multi-agent setting?
  \end{tcolorbox}
\end{frame}



	
	
	
	
\end{document}